%=======================================================
\chapter{Verbatim}

\textsc{Qual o pacote?}: \textbf{verbatim}\footnote{Por padrão, já incluído no \LaTeX.} (Acesso: \url{https://www.ctan.org/pkg/verbatim})

\textsc{O que faz?}: Digita qualquer coisa fora da formatação do texto, utilizando fonte monoespaçada e conforme os espaçamentos desejados. Útil para pequenos trechos de código.
\ \\

\begin{codex}{Códigos ambiente verbatim}
	\verb|comando 1|
	\verb*|comando 2]|
	\begin{verbatim}
		qualquer outro texto
	\end{verbatim}
\end{codex}
\ \\
\textsc{Resultado}:

\verb|comando 1|

\verb*|comando 2]|

\begin{verbatim}
	qualquer outro texto
	e outro texto qualquer
\end{verbatim}

\textbf{Importante:} O ambiente verbatim não respeitas as margens/quebra de linhas. Isto quer dizer que se tudo for digitado em uma única linha, sem quebra manual, ele irá mostrar o texto quebrado. Observe:

\begin{verbatim}
	Stricto sensu é uma expressão latina que significa, literalmente, "em sentido específico", por oposição ao "sentido amplo" de um termo. No âmbito do ensino, se refere ao nível de pós-graduação que titula o estudante como mestre ou doutor em determinado campo do conhecimento. Denota, neste caso, um grau mais elevado do que a pesquisa lato sensu.
\end{verbatim}

O código acima é:

\begin{codex}{Código verbatim sem quebra de linha}
Stricto sensu é uma expressão latina que significa, literalmente, "em sentido específico", por oposição ao "sentido amplo" de um termo. No âmbito do ensino, se refere ao nível de pós-graduação que titula o estudante como mestre ou doutor em determinado campo do conhecimento. Denota, neste caso, um grau mais elevado do que a pesquisa lato sensu.
\end{codex}